% Options for packages loaded elsewhere
\PassOptionsToPackage{unicode}{hyperref}
\PassOptionsToPackage{hyphens}{url}
\PassOptionsToPackage{dvipsnames,svgnames,x11names}{xcolor}
\documentclass[
  11pt,
  a4paper,
]{article}
\usepackage{xcolor}
\usepackage[margin=1in]{geometry}
\usepackage{amsmath,amssymb}
\setcounter{secnumdepth}{-\maxdimen} % remove section numbering
\usepackage{iftex}
\ifPDFTeX
  \usepackage[T1]{fontenc}
  \usepackage[utf8]{inputenc}
  \usepackage{textcomp} % provide euro and other symbols
\else % if luatex or xetex
  \usepackage{unicode-math} % this also loads fontspec
  \defaultfontfeatures{Scale=MatchLowercase}
  \defaultfontfeatures[\rmfamily]{Ligatures=TeX,Scale=1}
\fi
\usepackage{lmodern}
\ifPDFTeX\else
  % xetex/luatex font selection
  \ifLuaTeX
    \usepackage{luaotfload}
    \directlua{luaotfload.add_fallback("mainfontfallback",{
      "NotoSerifBengali:mode=harf"
    })}
  \fi
  \setmainfont[,RawFeature={fallback=mainfontfallback}]{Latin Modern Roman}
  \ifLuaTeX
    \usepackage{luaotfload}
    \directlua{luaotfload.add_fallback("sansfontfallback",{
      "NotoSansBengali:mode=harf"
    })}
  \fi
  \setsansfont[,RawFeature={fallback=sansfontfallback}]{Latin Modern Sans}
  \ifLuaTeX
    \usepackage{luaotfload}
    \directlua{luaotfload.add_fallback("monofontfallback",{
      "NotoSansBengali:mode=harf"
    })}
  \fi
  \setmonofont[,RawFeature={fallback=monofontfallback}]{DejaVu Sans Mono}
\fi
% Use upquote if available, for straight quotes in verbatim environments
\IfFileExists{upquote.sty}{\usepackage{upquote}}{}
\IfFileExists{microtype.sty}{% use microtype if available
  \usepackage[]{microtype}
  \UseMicrotypeSet[protrusion]{basicmath} % disable protrusion for tt fonts
}{}
\makeatletter
\@ifundefined{KOMAClassName}{% if non-KOMA class
  \IfFileExists{parskip.sty}{%
    \usepackage{parskip}
  }{% else
    \setlength{\parindent}{0pt}
    \setlength{\parskip}{6pt plus 2pt minus 1pt}}
}{% if KOMA class
  \KOMAoptions{parskip=half}}
\makeatother
\usepackage{color}
\usepackage{fancyvrb}
\newcommand{\VerbBar}{|}
\newcommand{\VERB}{\Verb[commandchars=\\\{\}]}
\DefineVerbatimEnvironment{Highlighting}{Verbatim}{commandchars=\\\{\}}
% Add ',fontsize=\small' for more characters per line
\usepackage{framed}
\definecolor{shadecolor}{RGB}{248,248,248}
\newenvironment{Shaded}{\begin{snugshade}}{\end{snugshade}}
\newcommand{\AlertTok}[1]{\textcolor[rgb]{0.94,0.16,0.16}{#1}}
\newcommand{\AnnotationTok}[1]{\textcolor[rgb]{0.56,0.35,0.01}{\textbf{\textit{#1}}}}
\newcommand{\AttributeTok}[1]{\textcolor[rgb]{0.13,0.29,0.53}{#1}}
\newcommand{\BaseNTok}[1]{\textcolor[rgb]{0.00,0.00,0.81}{#1}}
\newcommand{\BuiltInTok}[1]{#1}
\newcommand{\CharTok}[1]{\textcolor[rgb]{0.31,0.60,0.02}{#1}}
\newcommand{\CommentTok}[1]{\textcolor[rgb]{0.56,0.35,0.01}{\textit{#1}}}
\newcommand{\CommentVarTok}[1]{\textcolor[rgb]{0.56,0.35,0.01}{\textbf{\textit{#1}}}}
\newcommand{\ConstantTok}[1]{\textcolor[rgb]{0.56,0.35,0.01}{#1}}
\newcommand{\ControlFlowTok}[1]{\textcolor[rgb]{0.13,0.29,0.53}{\textbf{#1}}}
\newcommand{\DataTypeTok}[1]{\textcolor[rgb]{0.13,0.29,0.53}{#1}}
\newcommand{\DecValTok}[1]{\textcolor[rgb]{0.00,0.00,0.81}{#1}}
\newcommand{\DocumentationTok}[1]{\textcolor[rgb]{0.56,0.35,0.01}{\textbf{\textit{#1}}}}
\newcommand{\ErrorTok}[1]{\textcolor[rgb]{0.64,0.00,0.00}{\textbf{#1}}}
\newcommand{\ExtensionTok}[1]{#1}
\newcommand{\FloatTok}[1]{\textcolor[rgb]{0.00,0.00,0.81}{#1}}
\newcommand{\FunctionTok}[1]{\textcolor[rgb]{0.13,0.29,0.53}{\textbf{#1}}}
\newcommand{\ImportTok}[1]{#1}
\newcommand{\InformationTok}[1]{\textcolor[rgb]{0.56,0.35,0.01}{\textbf{\textit{#1}}}}
\newcommand{\KeywordTok}[1]{\textcolor[rgb]{0.13,0.29,0.53}{\textbf{#1}}}
\newcommand{\NormalTok}[1]{#1}
\newcommand{\OperatorTok}[1]{\textcolor[rgb]{0.81,0.36,0.00}{\textbf{#1}}}
\newcommand{\OtherTok}[1]{\textcolor[rgb]{0.56,0.35,0.01}{#1}}
\newcommand{\PreprocessorTok}[1]{\textcolor[rgb]{0.56,0.35,0.01}{\textit{#1}}}
\newcommand{\RegionMarkerTok}[1]{#1}
\newcommand{\SpecialCharTok}[1]{\textcolor[rgb]{0.81,0.36,0.00}{\textbf{#1}}}
\newcommand{\SpecialStringTok}[1]{\textcolor[rgb]{0.31,0.60,0.02}{#1}}
\newcommand{\StringTok}[1]{\textcolor[rgb]{0.31,0.60,0.02}{#1}}
\newcommand{\VariableTok}[1]{\textcolor[rgb]{0.00,0.00,0.00}{#1}}
\newcommand{\VerbatimStringTok}[1]{\textcolor[rgb]{0.31,0.60,0.02}{#1}}
\newcommand{\WarningTok}[1]{\textcolor[rgb]{0.56,0.35,0.01}{\textbf{\textit{#1}}}}
\usepackage{longtable,booktabs,array}
\usepackage{calc} % for calculating minipage widths
% Correct order of tables after \paragraph or \subparagraph
\usepackage{etoolbox}
\makeatletter
\patchcmd\longtable{\par}{\if@noskipsec\mbox{}\fi\par}{}{}
\makeatother
% Allow footnotes in longtable head/foot
\IfFileExists{footnotehyper.sty}{\usepackage{footnotehyper}}{\usepackage{footnote}}
\makesavenoteenv{longtable}
\setlength{\emergencystretch}{3em} % prevent overfull lines
\providecommand{\tightlist}{%
  \setlength{\itemsep}{0pt}\setlength{\parskip}{0pt}}
% paper_preamble.tex — shared LaTeX preamble for the PCMA/SI paper set.
%
% Compiled with lualatex (NOT pdflatex) because several papers contain
% Bengali script: the CLL paper realises one cognitive state into
% English, Bengali and Banglish, so the Bengali output is EVIDENCE and
% cannot be transliterated away without destroying the point being made.
%
% Font selection is NOT done here. pandoc >= 3 has native
% `mainfontfallback` / `monofontfallback` variables that emit the
% luaotfload fallback registration in the correct position relative to
% unicode-math; doing it by hand in this header lands after unicode-math
% and breaks font loading outright. See tools/build_papers.sh.

\usepackage{amsmath}
\usepackage{amssymb}
\usepackage{mathtools}   % \xleftrightarrow etc. used in cll.md §5
\usepackage{booktabs}

% ---------------------------------------------------------------------
% Unicode symbols appearing in PROSE rather than inside $...$.
%
% The papers were written as markdown and read as markdown, so they use
% real characters — "the ω-boost", "α-child", "ΔR coherence" — instead of
% LaTeX macros. Under unicode-math those codepoints are math-mode-only
% and abort the build ("\symbballowed only in math mode"). Mapping each
% to \ensuremath keeps the markdown readable on GitHub AND compilable,
% which is the whole point of keeping one source.
% ---------------------------------------------------------------------
\usepackage{newunicodechar}

% Greek, lower case
\newunicodechar{α}{\ensuremath{\alpha}}
\newunicodechar{β}{\ensuremath{\beta}}
\newunicodechar{γ}{\ensuremath{\gamma}}
\newunicodechar{δ}{\ensuremath{\delta}}
\newunicodechar{ε}{\ensuremath{\varepsilon}}
\newunicodechar{η}{\ensuremath{\eta}}
\newunicodechar{θ}{\ensuremath{\theta}}
\newunicodechar{κ}{\ensuremath{\kappa}}
\newunicodechar{λ}{\ensuremath{\lambda}}
\newunicodechar{μ}{\ensuremath{\mu}}
\newunicodechar{π}{\ensuremath{\pi}}
\newunicodechar{ρ}{\ensuremath{\rho}}
\newunicodechar{σ}{\ensuremath{\sigma}}
\newunicodechar{τ}{\ensuremath{\tau}}
\newunicodechar{φ}{\ensuremath{\varphi}}
\newunicodechar{χ}{\ensuremath{\chi}}
\newunicodechar{ψ}{\ensuremath{\psi}}
\newunicodechar{ω}{\ensuremath{\omega}}
\newunicodechar{ν}{\ensuremath{\nu}}
\newunicodechar{ξ}{\ensuremath{\xi}}
\newunicodechar{ζ}{\ensuremath{\zeta}}
\newunicodechar{ι}{\ensuremath{\iota}}

% Greek, upper case
\newunicodechar{Δ}{\ensuremath{\Delta}}
\newunicodechar{Φ}{\ensuremath{\Phi}}
\newunicodechar{Σ}{\ensuremath{\Sigma}}
\newunicodechar{Ψ}{\ensuremath{\Psi}}
\newunicodechar{Ω}{\ensuremath{\Omega}}
\newunicodechar{Γ}{\ensuremath{\Gamma}}
\newunicodechar{Λ}{\ensuremath{\Lambda}}
\newunicodechar{Θ}{\ensuremath{\Theta}}
\newunicodechar{Π}{\ensuremath{\Pi}}

% Operators and relations
\newunicodechar{×}{\ensuremath{\times}}
\newunicodechar{·}{\ensuremath{\cdot}}
\newunicodechar{±}{\ensuremath{\pm}}
\newunicodechar{−}{\ensuremath{-}}
\newunicodechar{≈}{\ensuremath{\approx}}
\newunicodechar{≤}{\ensuremath{\leq}}
\newunicodechar{≥}{\ensuremath{\geq}}
\newunicodechar{≠}{\ensuremath{\neq}}
\newunicodechar{∈}{\ensuremath{\in}}
\newunicodechar{∂}{\ensuremath{\partial}}
\newunicodechar{∇}{\ensuremath{\nabla}}
\newunicodechar{√}{\ensuremath{\sqrt{}}}
\newunicodechar{∞}{\ensuremath{\infty}}
\newunicodechar{∝}{\ensuremath{\propto}}
\newunicodechar{∑}{\ensuremath{\sum}}
\newunicodechar{∫}{\ensuremath{\int}}
\newunicodechar{ℝ}{\ensuremath{\mathbb{R}}}
\newunicodechar{ℕ}{\ensuremath{\mathbb{N}}}
\newunicodechar{‖}{\ensuremath{\|}}
\newunicodechar{⊕}{\ensuremath{\oplus}}
\newunicodechar{⊗}{\ensuremath{\otimes}}
\newunicodechar{∀}{\ensuremath{\forall}}
\newunicodechar{∃}{\ensuremath{\exists}}
\newunicodechar{¬}{\ensuremath{\neg}}
\newunicodechar{∧}{\ensuremath{\wedge}}
\newunicodechar{∨}{\ensuremath{\vee}}

% Arrows
\newunicodechar{→}{\ensuremath{\rightarrow}}
\newunicodechar{←}{\ensuremath{\leftarrow}}
\newunicodechar{↔}{\ensuremath{\leftrightarrow}}
\newunicodechar{↑}{\ensuremath{\uparrow}}
\newunicodechar{↓}{\ensuremath{\downarrow}}
\newunicodechar{⇒}{\ensuremath{\Rightarrow}}
\newunicodechar{⇐}{\ensuremath{\Leftarrow}}
\newunicodechar{⇔}{\ensuremath{\Leftrightarrow}}

% Super/subscripts written as literal characters
\newunicodechar{²}{\ensuremath{^{2}}}
\newunicodechar{³}{\ensuremath{^{3}}}
\newunicodechar{⁴}{\ensuremath{^{4}}}
\newunicodechar{⁵}{\ensuremath{^{5}}}
\newunicodechar{⁶}{\ensuremath{^{6}}}
\newunicodechar{₀}{\ensuremath{_{0}}}
\newunicodechar{₁}{\ensuremath{_{1}}}
\newunicodechar{₂}{\ensuremath{_{2}}}
\newunicodechar{₃}{\ensuremath{_{3}}}
\newunicodechar{ᵀ}{\ensuremath{^{\mathsf{T}}}}

% Status glyphs. The papers use emoji as status markers in tables
% (implemented / failed / warning). Mapped to typographic equivalents:
% a colour emoji font in lualatex is unreliable, and the MEANING is what
% carries — a reader needs to see "this one failed", not a red cross.
\newunicodechar{✅}{\ensuremath{\checkmark}}
\newunicodechar{✓}{\ensuremath{\checkmark}}
\newunicodechar{❌}{\ensuremath{\times}}
\newunicodechar{✗}{\ensuremath{\times}}
\newunicodechar{⚠}{\textbf{!}}
\newunicodechar{🧩}{\textbf{[piece]}}
\newunicodechar{📋}{\textbf{[list]}}
\newunicodechar{🟡}{\textbf{[partial]}}
\newunicodechar{🔴}{\textbf{[dead]}}
\newunicodechar{🟢}{\textbf{[live]}}
\newunicodechar{▼}{\ensuremath{\blacktriangledown}}
\newunicodechar{►}{\ensuremath{\blacktriangleright}}
\newunicodechar{□}{\ensuremath{\square}}
\newunicodechar{█}{\rule{0.6em}{0.9ex}}
\newunicodechar{•}{\textbullet}
\newunicodechar{†}{\textdagger}
\newunicodechar{‑}{-}
\newunicodechar{‐}{-}
\newunicodechar{…}{\ldots}
\newunicodechar{।}{\textbar}   % danda: rendered by the Bengali fallback in text
\newunicodechar{⁻}{\ensuremath{^{-}}}

% Zero-width formatting codepoints. These are invisible in an editor but
% are real characters to LaTeX: an emoji written as "⚠️" is TWO
% codepoints (U+26A0 U+FE0F), so declaring only the visible one leaves a
% variation selector with no definition and the build dies on a symbol
% nobody can see. Both are mapped to nothing.
\newunicodechar{^^^^fe0f}{}    % VARIATION SELECTOR-16
\newunicodechar{^^^^200d}{}    % ZERO WIDTH JOINER
\newunicodechar{^^^^200b}{}    % ZERO WIDTH SPACE

\usepackage{longtable}
\usepackage{array}

% Comparison matrices and cost ledgers are wide. Let them shrink to the
% text block rather than overflow the margin.
\usepackage{adjustbox}

% Code blocks: wrap long lines instead of running off the page. Several
% papers quote real C++ with long identifiers.
\usepackage{fvextra}
\fvset{breaklines=true,breakanywhere=true,fontsize=\small}

\setlength{\parskip}{0.5em}
\setlength{\parindent}{0pt}

\usepackage{fancyhdr}
\pagestyle{fancy}
\fancyhf{}
\fancyfoot[C]{\thepage}
\fancyhead[L]{\small\itshape Synthetic Intelligence — research lineage}
\fancyhead[R]{\small\itshape irfan@furylogic.com}
\renewcommand{\headrulewidth}{0.2pt}
\setlength{\headheight}{14pt}
\setlength{\headsep}{18pt}
\usepackage{bookmark}
\IfFileExists{xurl.sty}{\usepackage{xurl}}{} % add URL line breaks if available
\urlstyle{same}
\hypersetup{
  pdftitle={DCHL Thinking Machine},
  pdfauthor={Irfan Mahir --- Independent Researcher, Dhaka, Bangladesh --- irfan@furylogic.com},
  colorlinks=true,
  linkcolor={black},
  filecolor={Maroon},
  citecolor={Blue},
  urlcolor={[HTML]{0B5394}},
  pdfcreator={LaTeX via pandoc}}

\title{DCHL Thinking Machine}
\author{Irfan Mahir --- Independent Researcher, Dhaka, Bangladesh ---
irfan@furylogic.com}
\date{}

\begin{document}
\maketitle

{
\hypersetup{linkcolor=}
\setcounter{tocdepth}{3}
\tableofcontents
}
\begin{quote}
\textbf{STATUS: RETIRED ARCHITECTURE.} This folder documents the DCHL
substrate, which was replaced by the oscillator substrate. It is kept
because the negative result is a contribution --- SI paper §2 is a
self-critique built on it. Only RIFA (IC/SIL) remains load-bearing. See
SI paper §6.2.2. Papers in \texttt{DEAD/} are formally dead: §6.2.3.
\end{quote}

\section{DCHL Thinking Machine}\label{dchl-thinking-machine}

\textbf{Irfan Mahir} \emph{Independent Researcher, Dhaka, Bangladesh}
irfan@furylogic.com

\textbf{Decoupled Contrastive Hebbian Learning} --- a physics-based
cognitive architecture that \emph{thinks} rather than \emph{generates}.
No tokens, no attention, no backprop, no GPU needed.

Input → byte-mapped cognitive state (256 floats) → Hamiltonian reasoning
on learned potential landscape → deeper understanding (lower Φ).

\subsection{Quick Start}\label{quick-start}

\begin{Shaded}
\begin{Highlighting}[]
\FunctionTok{make} \AttributeTok{{-}B}\NormalTok{ chatbot lto      }\CommentTok{\# build (CPU only, g++ 9+)}
\ExtensionTok{./lto} \AttributeTok{{-}{-}epochs}\NormalTok{ 200 }\AttributeTok{{-}{-}save}\NormalTok{ chatbot.neuron  }\CommentTok{\# train \textasciitilde{}3 min}
\ExtensionTok{./chatbot}\NormalTok{ chatbot.neuron                    }\CommentTok{\# think interactively}
\end{Highlighting}
\end{Shaded}

\begin{verbatim}
=== DCHL Thinking Machine ===
Input → cognitive state → Hamiltonian reasoning
The state IS the understanding.

> what is the capital of france?
  Φ(initial)=-295.1  conf=0.999  ood=0
  Φ(final)=-737.3    ΔΦ=-442.2  depth=100
  ✓ Deepened understanding
\end{verbatim}

\subsection{What It Is}\label{what-it-is}

A \textbf{thinker}, not a generator. The model compresses text directly
into a 256-float cognitive state (byte→float, no tokenizer), then
evolves that state through a learned energy landscape via Hamiltonian
dynamics. The evolved state IS the understanding --- not a text string,
not a probability distribution over tokens, but a compressed cognitive
representation.

\subsubsection{What It Can Do (Same Tasks as ML/DL --- Cheaper, CPU
Only)}\label{what-it-can-do-same-tasks-as-mldl-cheaper-cpu-only}

\begin{longtable}[]{@{}
  >{\raggedright\arraybackslash}p{(\linewidth - 6\tabcolsep) * \real{0.2069}}
  >{\raggedright\arraybackslash}p{(\linewidth - 6\tabcolsep) * \real{0.1724}}
  >{\raggedright\arraybackslash}p{(\linewidth - 6\tabcolsep) * \real{0.3448}}
  >{\raggedright\arraybackslash}p{(\linewidth - 6\tabcolsep) * \real{0.2759}}@{}}
\toprule\noalign{}
\begin{minipage}[b]{\linewidth}\raggedright
Task
\end{minipage} & \begin{minipage}[b]{\linewidth}\raggedright
How
\end{minipage} & \begin{minipage}[b]{\linewidth}\raggedright
vs ML/DL
\end{minipage} & \begin{minipage}[b]{\linewidth}\raggedright
vs LLM
\end{minipage} \\
\midrule\noalign{}
\endhead
\bottomrule\noalign{}
\endlastfoot
Classification & Lightweight head on evolved state & No GPU needed & No
GPU needed \\
Similarity & Cosine between cognitive states & No embedding model & No
embedding model \\
Anomaly & Φ(c) IS the anomaly score --- built-in & No OOD training & No
OOD training \\
Retrieval & State-to-state nearest neighbor & No index needed & No
KV-cache needed \\
Reasoning & Hamiltonian dynamics → deeper understanding & No
chain-of-thought & No LLM inference \\
\end{longtable}

\subsubsection{What It Is NOT}\label{what-it-is-not}

\begin{itemize}
\tightlist
\item
  ❌ Not a text generator (doesn't produce output tokens)
\item
  ❌ Not a transformer (no attention, no KV-cache, no context window)
\item
  ❌ Not an LLM (33K params, not billions)
\item
  ❌ Not gradient-based (no backprop, no loss functions, no Adam)
\item
  ❌ Not a statistical pattern matcher (Hamiltonian mechanics on learned
  energy landscape)
\end{itemize}

\subsection{Architecture}\label{architecture}

\begin{verbatim}
State (256 floats) → Hidden (128, tanh) → Output (1 float = Φ)
         ↑                                      ↑
    byte→float mapping               energy of cognitive state
    (no tokenizer)                   (lower = more "understood")
\end{verbatim}

Two independent learning processes, both local, no backprop: -
\textbf{Hidden layer}: per-unit contrastive Hebbian --- units
differentiate text from random - \textbf{Output layer}: pure
differential projection --- maps features to energy values

Training is text-vs-random: Φ(text) \textless{} Φ(random), that's all.
No transitions, no predictions, no targets.

\subsubsection{Training vs Inference}\label{training-vs-inference}

\begin{verbatim}
Training:
  For each pair (text, random):
    Φ(text)  → push lower (h_pos - h_neg)
    Φ(random) → push higher (h_neg - h_pos)

Inference (Hamiltonian reasoning):
  c₀ = byte→float(input)
  for step in 1..100:
    p += -∇Φ(c) · dt      (gradient of potential = force)
    c += p · dt            (momentum drives state)
  c_final = evolved, more "understood" state
  ΔΦ = Φ(c_final) - Φ(c₀)  (negative = understanding deepened)
\end{verbatim}

\subsection{Hardware}\label{hardware}

\textbf{Runs on commodity CPU servers. The same old hardware used for
spam detection.}

\begin{itemize}
\tightlist
\item
  No GPU required
\item
  \textasciitilde33K parameters (fits in L2 cache)
\item
  Training: \textasciitilde3 minutes for 200 epochs
\item
  Inference: \textasciitilde0.01s for 100 reasoning steps
\item
  Model file: \textasciitilde264 KB
\end{itemize}

\subsubsection{vs The Industry}\label{vs-the-industry}

\begin{longtable}[]{@{}lllll@{}}
\toprule\noalign{}
Architecture & Params & Hardware & Training & Inference \\
\midrule\noalign{}
\endhead
\bottomrule\noalign{}
\endlastfoot
DCHL (this) & 33K & Any CPU & \textasciitilde3 min &
\textasciitilde0.01s \\
BERT-base & 110M & GPU & days & \textasciitilde0.1s \\
GPT-3 & 175B & H100 cluster & months & \textasciitilde1s \\
Llama 3 70B & 70B & H100 cluster & weeks & \textasciitilde0.5s \\
\end{longtable}

DCHL inverts the scaling law: better physics → deeper understanding ---
same compute. Not bigger models → diminishing returns.

\subsection{Files}\label{files}

\begin{longtable}[]{@{}
  >{\raggedright\arraybackslash}p{(\linewidth - 2\tabcolsep) * \real{0.4000}}
  >{\raggedright\arraybackslash}p{(\linewidth - 2\tabcolsep) * \real{0.6000}}@{}}
\toprule\noalign{}
\begin{minipage}[b]{\linewidth}\raggedright
File
\end{minipage} & \begin{minipage}[b]{\linewidth}\raggedright
Purpose
\end{minipage} \\
\midrule\noalign{}
\endhead
\bottomrule\noalign{}
\endlastfoot
\texttt{src/hcma.hpp} & DCHL potential, contrastive Hebbian updates,
Hamiltonian dynamics \\
\texttt{src/lto.cpp} & Training entry point (CPU/GPU, CLI, eval) \\
\texttt{src/chatbot.cpp} & Interactive thinking machine \\
\texttt{src/bench.cpp} & Performance benchmarking \\
\texttt{src/gpu\_training.cu} & CUDA kernels (optional) \\
\end{longtable}

\subsection{CLI}\label{cli}

\subsubsection{\texorpdfstring{Training
(\texttt{lto})}{Training (lto)}}\label{training-lto}

\begin{Shaded}
\begin{Highlighting}[]
\ExtensionTok{./lto} \AttributeTok{{-}{-}epochs}\NormalTok{ 200 }\AttributeTok{{-}{-}lr{-}hidden}\NormalTok{ 0.01 }\AttributeTok{{-}{-}lr{-}output}\NormalTok{ 0.1 }\AttributeTok{{-}{-}save}\NormalTok{ model.neuron}
\end{Highlighting}
\end{Shaded}

Key flags: \texttt{-\/-epochs}, \texttt{-\/-lr-hidden},
\texttt{-\/-lr-output}, \texttt{-\/-d-state}, \texttt{-\/-d-hidden},
\texttt{-\/-data}

\subsubsection{\texorpdfstring{Chat
(\texttt{chatbot})}{Chat (chatbot)}}\label{chat-chatbot}

\begin{Shaded}
\begin{Highlighting}[]
\ExtensionTok{./chatbot}\NormalTok{ model.neuron }\AttributeTok{{-}{-}model}\NormalTok{ potential.neuron }\AttributeTok{{-}{-}steps}\NormalTok{ 100 }\AttributeTok{{-}{-}dt}\NormalTok{ 0.005}
\end{Highlighting}
\end{Shaded}

Shows Φ(initial), Φ(final), ΔΦ, confidence, OOD flag.

\subsection{Research}\label{research}

\begin{itemize}
\tightlist
\item
  \href{research-paper/hcma.md}{HCMA paper}: Full architecture \& math
\item
  \href{contributing/how_we_implement.md}{Implementation notes}: Every
  problem \& fix
\end{itemize}

\subsection{License}\label{license}

MIT

\end{document}
